\documentclass[11pt,a4paper]{article}
\usepackage[utf8]{inputenc}
\usepackage[T1]{fontenc}
\usepackage{amsmath,amssymb,amsthm}
\usepackage[margin=2.6cm]{geometry}
\usepackage{hyperref}
\usepackage{booktabs}
\usepackage{listings}
\lstset{basicstyle=\ttfamily\small,breaklines=true,columns=fullflexible,keepspaces=true,
  escapechar=@,frame=single,xleftmargin=4pt,xrightmargin=4pt}

\newtheorem{theorem}{Theorem}[section]
\newtheorem{lemma}[theorem]{Lemma}
\newtheorem{proposition}[theorem]{Proposition}
\newtheorem{corollary}[theorem]{Corollary}
\theoremstyle{definition}
\newtheorem{definition}[theorem]{Definition}
\theoremstyle{remark}
\newtheorem{remark}[theorem]{Remark}

\newcommand{\F}{\mathbb{F}}
\newcommand{\Lset}{\mathcal{L}}
\newcommand{\Oset}{\mathcal{O}}

\title{Conjecture 00000001064: the minimal Kakeya set in $\F_q^2$ for odd $q$\\
\large A proof, with a Lean~4 / Mathlib formalization}
\author{Zhiyuan Zhang (GitHub: \texttt{LibertychaserUS})}
\date{27 October 2026}

\begin{document}
\maketitle

\begin{abstract}
We prove conjecture 00000001064 of \emph{The Last Math Competition}. Let $\F$ be a finite field
with $q=|\F|$ odd. Every Kakeya set in the affine plane $\F^2$ (a set containing a full affine line
in each of the $q+1$ directions) has at least $q(q+1)/2+(q-1)/2$ points, and this bound is attained
by the explicit set $\{(x,y): x^2-y \text{ is a square}\}\cup\{x=0\}$. The lower bound is the
theorem of Blokhuis and Mazzocca (2008). The argument combines Faber's incidence identity for line
systems, a parity argument producing a blocking set from the points of odd multiplicity, and
Jamison's theorem on blocking sets of $AG(2,q)$, which is proved here by the polynomial method.
The whole statement is formalized in Lean~4 with Mathlib (Lean~4 toolchain and Mathlib \texttt{v4.33.1})
in the file \texttt{Kakeya.lean}; it compiles without errors or warnings, and the eight main
declarations depend only on \texttt{propext}, \texttt{Classical.choice} and \texttt{Quot.sound}.
\end{abstract}

\tableofcontents

\section{The conjecture as filed}

The conjecture is filed as follows (English version; the Chinese version in the submitted
\texttt{.md} file is a literal translation of this text).

\begin{quote}
\textbf{Conjecture 00000001064.} The minimal Kakeya set in $\F_q^2$ for odd $q$ has size
$\dfrac{q(q+1)}{2}+\left\lfloor\dfrac{q-1}{2}\right\rfloor$ (the exact minimum for odd $q$).
\end{quote}

\section{Precise reading and the theorem proved}

Throughout, $\F$ is a finite field with $q=|\F|$ elements, and the affine plane is $\F^2$ (the
affine plane $AG(2,q)$). The following conventions make the statement precise.

\paragraph{Lines and directions.} For $m,b\in\F$ put
$L(m,b)=\{(x,\,mx+b): x\in\F\}$, the non-vertical affine line of slope $m$ through $(0,b)$, and for
$c\in\F$ put $V(c)=\{(c,y): y\in\F\}$, the vertical line $x=c$. Each line has exactly $q$ points.
The \emph{directions} are the $q+1$ elements of $\F\cup\{\infty\}$: the direction $m\in\F$ consists
of all lines $L(m,b)$, $b\in\F$, and the direction $\infty$ consists of all lines $V(c)$, $c\in\F$.
Two lines in the same direction are parallel: they are equal or disjoint. Two lines in different
directions meet in exactly one point.

\paragraph{Kakeya sets.} A \emph{Kakeya set} is a set $K\subseteq\F^2$ which contains a full
affine line in every direction. Explicitly, $K$ is Kakeya if for every $m\in\F$ there is $b\in\F$
with $L(m,b)\subseteq K$, and there is $c\in\F$ with $V(c)\subseteq K$. The conjecture concerns the
minimum $\mu(q)=\min\{|K|: K \text{ Kakeya}\}$.

\paragraph{The floor.} For odd $q$ the integer $(q-1)/2$ is an integer, so
$\lfloor (q-1)/2\rfloor=(q-1)/2$, and the claimed value is
$\dfrac{q(q+1)}{2}+\dfrac{q-1}{2}=\dfrac{q^2+2q-1}{2}$. For example $q=3,5,7,11,13$ give $7,17,31,71,97$.

\begin{theorem}[Blokhuis--Mazzocca, conjecture 00000001064]\label{thm:main}
Let $\F$ be a finite field with $q=|\F|$ odd.
\begin{enumerate}
\item[(a)] Every Kakeya set $K\subseteq\F^2$ satisfies $|K|\ge \dfrac{q(q+1)}{2}+\dfrac{q-1}{2}$.
\item[(b)] The set $E=\{(x,y)\in\F^2 : x^2-y\text{ is a square in }\F\}\cup\{(0,y): y\in\F\}$ is a
Kakeya set with $|E|=\dfrac{q(q+1)}{2}+\dfrac{q-1}{2}$.
\end{enumerate}
Consequently the minimum size of a Kakeya set in $\F^2$ is exactly $\dfrac{q(q+1)}{2}+\dfrac{q-1}{2}$.
\end{theorem}

\paragraph{History.} Faber stated the planar value as a conjecture. Wolff's approach gave an
earlier bound of order $q^2/2$ for the planar problem, and the incidence count that yields
$q(q+1)/2$ for every $q$ is credited by Blokhuis and Mazzocca to Faber's counting argument. Dvir
(2009) proved that Kakeya sets in $\F_q^n$ have size at least $c_n q^n$ with $c_n$ depending only on
$n$; this is of the right order but not sharp in the plane. Blokhuis and Mazzocca (2008) settled
Faber's conjecture for odd $q$ and characterized the extremal example. The even case is not
treated here. The proof below is reconstructed from the affine-coordinate argument in the
accompanying file \texttt{PROOF.md}; the only deviation from the paper is that the product
computation of Section~3.4 replaces the appeal to Bichara and Korchm\'aros, and Jamison's theorem is
proved in full.

\section{Mathematical proof}

Fix $\F$ with $q$ odd, and fix a Kakeya set $K$ unless stated otherwise. The proof of the lower
bound has three steps: reduction to a line system (Section~3.1), the easy bound (Section~3.2), and
the refinement for odd $q$ (Sections~3.3 and 3.4). Section~3.5 gives the construction.

\subsection{Line systems and Faber's incidence identity}

\paragraph{Line systems.} If $K$ is Kakeya, choose for every direction $d\in\F\cup\{\infty\}$ a
line $L_d\subseteq K$ of direction $d$. Write $L_m=L(m,t_m)$ and $L_\infty=V(t_\infty)$; the family
$\Lset=\{L_d\}$ is a \emph{line system}, with exactly one line in each of the $q+1$ directions. Its
union $U(\Lset)=\bigcup_d L_d$ satisfies $U(\Lset)\subseteq K$, so $|K|\ge|U(\Lset)|$. It therefore
suffices to bound $|U(\Lset)|$ from below for every line system.

For a point $A\in\F^2$ let $n(A)=\#\{d: A\in L_d\}$ be its \emph{multiplicity}. Then
$U(\Lset)=\{A: n(A)\ge 1\}$.

\begin{lemma}[basic counts]\label{lem:counts}
For every line system $\Lset$:
\begin{enumerate}
\item[(C1)] $\sum_{A} n(A)=\sum_d|L_d|=q(q+1)$.
\item[(C2)] $\sum_A n(A)(n(A)-1)=\#\{(d,d',A): d\ne d',\ A\in L_d\cap L_{d'}\}=q(q+1)$.
\end{enumerate}
Consequently $\sum_A n(A)(n(A)-2)=0$.
\end{lemma}

\begin{proof}
(C1) is the double count of pairs (direction, point of its line). For (C2), any two lines of
distinct directions meet in exactly one point, and the number of ordered pairs of distinct
directions is $(q+1)q$. The consequence is (C2) minus (C1).
\end{proof}

\begin{lemma}[Faber's identity]\label{lem:faber}
Let $\sigma=\tfrac12\sum_{A:\,n(A)\ge1}(n(A)-1)(n(A)-2)$. Then
\[
2\,|U(\Lset)| \;=\; q(q+1)+2\sigma .
\]
Moreover $\sigma\ge 0$.
\end{lemma}

\begin{proof}
For every integer $n\ge0$ one has $2[n\ge1]=2n-n(n-1)+[n\ge1](n-1)(n-2)$: for $n=0$ both sides are
$0$, and for $n\ge1$ the right side is $2n-n^2+n+n^2-3n+2=2$. Summing over $A\in\F^2$ and using (C1)
and (C2) gives $2|U(\Lset)|=2q(q+1)-q(q+1)+\sum_{n(A)\ge1}(n(A)-1)(n(A)-2)=q(q+1)+2\sigma$. Each
term $(n-1)(n-2)$ with $n\ge1$ is nonnegative, since $n-1\ge0$ and either $n-2\ge0$ or $n=1$.
\end{proof}

\subsection{The easy bound $q(q+1)/2$}

\begin{proposition}\label{prop:easy}
For every finite field $\F$ and every Kakeya set $K$, $|K|\ge q(q+1)/2$.
\end{proposition}

\begin{proof}
By Lemma~\ref{lem:faber} and $\sigma\ge0$ we get $|U(\Lset)|\ge q(q+1)/2$ for the line system
extracted from $K$, and $|K|\ge|U(\Lset)|$.
\end{proof}

\subsection{Jamison's theorem: blocking sets of $AG(2,q)$}

A set $B\subseteq\F^2$ is a \emph{blocking set} if it meets every affine line of $\F^2$ (all $q^2+q$
of them).

\begin{theorem}[Jamison]\label{thm:jamison}
Every blocking set $B\subseteq\F^2$ satisfies $|B|\ge 2q-1$.
\end{theorem}

\begin{proof}[Proof (Brouwer--Schrijver, dualised)]
Translating $B$ we may assume $(0,0)\in B$; translation preserves blocking sets. Put
\[
f(X,Z)=\prod_{(a,b)\in B,\ (a,b)\ne(0,0)}\bigl(1-aX+bZ\bigr)\in\F[X,Z].
\]
Then $\deg f\le |B|-1$, $f(0,0)=1$, and $f(x,z)=0$ for every $(x,z)\ne(0,0)$. Indeed, if $z\ne0$, the
line $bz=xa-1$ (equivalently $b=(x/z)a-1/z$) is non-vertical, so $B$ contains a point $(a,b)$ on it;
this point is not $(0,0)$ because $-1/z\ne0$, and $1-xa+bz=0$. If $z=0$ then $x\ne0$, the vertical
line $a=1/x$ meets $B$ in some $(1/x,b)$ with $(1/x,b)\ne(0,0)$, and the factor
$1-(1/x)x+b\cdot0$ vanishes.

Hence $\sum_{(x,z)\in\F^2}f(x,z)=f(0,0)=1$. On the other hand write $f=\sum_e c_e X^{e_1}Z^{e_2}$.
Then
\[
\sum_{(x,z)}f(x,z)=\sum_e c_e\Bigl(\sum_{x\in\F}x^{e_1}\Bigr)\Bigl(\sum_{z\in\F}z^{e_2}\Bigr),
\]
and $\sum_{x\in\F}x^{e}=0$ whenever $e<q-1$ (for $e=0$ the sum is $q=0$ in $\F$). If $\deg f<2(q-1)$,
then every monomial has $e_1<q-1$ or $e_2<q-1$, so the total sum is $0\ne1$. Therefore
$|B|-1\ge\deg f\ge2(q-1)$, i.e. $|B|\ge2q-1$.
\end{proof}

\subsection{The parity argument for odd $q$}

Let $\Lset$ be a line system with $q$ odd. Let $\Oset=\{A: n(A)\text{ odd}\}$ and $\tau=|\Oset|$.
Call a line $L_d\in\Lset$ \emph{all-2} if $n(A)=2$ for every $A\in L_d$.

\begin{lemma}[counting odd points]\label{lem:taubound}
$\tau\le 4\sigma$.
\end{lemma}

\begin{proof}
Pointwise, for every integer $n\ge0$,
\[
[n\text{ odd}]\;\le\;2\,[n\ge1](n-1)(n-2)-n(n-2).
\]
Check: $n=0$ gives $0\le0$; $n=1$ gives $1\le 1$; $n=2$ gives $0\le0$; $n=3$ gives $1\le 1$; and for
$n\ge4$ the right side equals $(n-2)^2\ge4$. Summing over $A$, and using $\sum_A n(A)(n(A)-2)=0$
from Lemma~\ref{lem:counts} and $\sum_{n(A)\ge1}(n(A)-1)(n(A)-2)=2\sigma$, gives $\tau\le 4\sigma$.
\end{proof}

\begin{lemma}[parity of lines]\label{lem:parity}
Let $q$ be odd.
\begin{enumerate}
\item Every affine line $L\subseteq\F^2$ with $L\notin\Lset$ contains a point of $\Oset$.
\item Every $L_d\in\Lset$ contains a point of $\Oset$ or is all-2.
\end{enumerate}
\end{lemma}

\begin{proof}
(1) $L$ has the direction of exactly one $L_d$, which is disjoint from $L$ (it is parallel and
different), while every other $L_{d'}$ meets $L$ in exactly one point. So $\sum_{A\in L}n(A)=q$,
which is odd, and some $A\in L$ has $n(A)$ odd.

(2) $\sum_{A\in L_d}n(A)=\sum_{A\in L_d}1+\sum_{A\in L_d}(n(A)-1)=q+q=2q$, because each of the other
$q$ lines meets $L_d$ in one point. Suppose $L_d$ contains no point of $\Oset$. Then every
$n(A)-1$ is odd and $\ge1$ for $A\in L_d$. These are $q$ odd positive integers with sum $q$, so each
equals $1$, i.e. $n(A)=2$ on $L_d$.
\end{proof}

\begin{lemma}[product lemma]\label{lem:product}
Let $q$ be odd. A line system $\Lset$ has at most one non-vertical all-2 line: if $L_{m_0}$ and
$L_{m_1}$ with $m_0\ne m_1$ are both all-2, we get a contradiction.
\end{lemma}

\begin{proof}
Let $P=(x^*,y^*)$ be the intersection point of $L_{m_0}$ and $L_{m_1}$, and let $c$ be the abscissa
of the vertical line $L_\infty=V(c)$. Write $b_m$ for the intercept of $L_m$. For $m\in S:=\F\setminus\{m_0,m_1\}$
define the abscissas of the intersections with $L_{m_0}$ and $L_{m_1}$:
\[
u(m)=\frac{b_{m_0}-b_m}{m-m_0},\qquad v(m)=\frac{b_{m_1}-b_m}{m-m_1},
\]
and $\gamma(m)=m x^*+b_m$. A direct computation, using $y^*=m_0x^*+b_{m_0}=m_1x^*+b_{m_1}$, gives
\[
u(m)-x^*=\frac{y^*-\gamma(m)}{m-m_0},\qquad v(m)-x^*=\frac{y^*-\gamma(m)}{m-m_1}.
\]
Since $L_{m_0}$ is all-2, each of its $q$ points lies on exactly one further line of $\Lset$. The
$q$ other lines of $\Lset$ meet $L_{m_0}$ in one point each, so these $q$ points are pairwise
distinct and exhaust $L_{m_0}$. Their abscissas are $u(m)$ for $m\in\F\setminus\{m_0\}$ together with
$c$ (for $L_\infty$). Hence $m\mapsto u(m)$ is a bijection $\F\setminus\{m_0\}\to\F\setminus\{c\}$, and
it sends $m_1$ to $x^*$. Restricting to $S$, $u$ is a bijection $S\to\F\setminus\{c,x^*\}$. The same
argument with $L_{m_1}$ shows that $v$ is a bijection $S\to\F\setminus\{c,x^*\}$. Therefore
\[
\prod_{m\in S}\bigl(u(m)-x^*\bigr)=\prod_{w\in\F\setminus\{c,x^*\}}(w-x^*)=\prod_{m\in S}\bigl(v(m)-x^*\bigr)=:R,
\]
and all factors are nonzero. Substituting the identities above,
\[
\frac{\prod_{S}(y^*-\gamma(m))}{\prod_{S}(m-m_0)}=R=\frac{\prod_{S}(y^*-\gamma(m))}{\prod_{S}(m-m_1)}.
\]
The numerator is nonzero, so cancelling gives $\prod_{S}(m-m_0)=\prod_{S}(m-m_1)=:P_0$, and $P_0\ne0$.
Let $W=\prod_{w\ne0}w$. Substituting $w=m-m_0$ (resp.\ $w=m-m_1$) gives
$(m_1-m_0)\,P_0=W$ and $(m_0-m_1)\,P_0=W$. Subtracting, $2(m_1-m_0)P_0=0$. This is impossible:
$2\ne0$ because $q$ is odd, $m_1\ne m_0$, and $P_0\ne0$.
\end{proof}

\begin{lemma}[blocking set]\label{lem:block}
Let $q$ be odd. There is a point $E\in L_\infty$ such that $E$ lies on every all-2 line of $\Lset$.
The set $B=\Oset\cup\{E\}$ is a blocking set, so $\tau+1\ge|B|\ge 2q-1$.
\end{lemma}

\begin{proof}
If some non-vertical $L_{m_0}$ is all-2, it is unique by Lemma~\ref{lem:product}; take
$E=L_\infty\cap L_{m_0}$. Otherwise take $E=(c,0)\in L_\infty$. In both cases $E$ lies on $L_\infty$
and on every all-2 line. Now let $L$ be any affine line. If $L\notin\Lset$, it meets $\Oset$ by
Lemma~\ref{lem:parity}(1). If $L\in\Lset$ is not all-2, it meets $\Oset$ by Lemma~\ref{lem:parity}(2).
If $L\in\Lset$ is all-2, then $L$ is $L_\infty$ or the unique non-vertical all-2 line, and in either
case $L\ni E$. So $B$ meets every affine line, and Theorem~\ref{thm:jamison} gives $|B|\ge2q-1$.
\end{proof}

\begin{proposition}\label{prop:sigma}
For $q$ odd and any line system, $\sigma\ge\dfrac{q-1}{2}$.
\end{proposition}

\begin{proof}
By Lemma~\ref{lem:block}, $\tau\ge 2q-2$. By Lemma~\ref{lem:taubound}, $4\sigma\ge\tau\ge2q-2$, so
$\sigma\ge(q-1)/2$.
\end{proof}

\begin{proof}[Proof of Theorem~\ref{thm:main}(a)]
For a Kakeya set $K$ take its line system $\Lset$ with $U(\Lset)\subseteq K$. By
Lemma~\ref{lem:faber} and Proposition~\ref{prop:sigma},
\[
|K|\ \ge\ |U(\Lset)|\ =\ \frac{q(q+1)}{2}+\sigma\ \ge\ \frac{q(q+1)}{2}+\frac{q-1}{2}.
\]
\end{proof}

\subsection{The construction}

Let $q$ be odd, so $2$ and $4$ are invertible in $\F$. Define
\[
E=\{(x,y)\in\F^2: x^2-y\text{ is a square in }\F\}\ \cup\ \{(0,y): y\in\F\}.
\]

\begin{proof}[Proof of Theorem~\ref{thm:main}(b)]
\emph{Kakeya.} For $m\in\F$ the line $y=mx-m^2/4$ lies in $E$, because
$x^2-(mx-m^2/4)=(x-m/2)^2$ is a square. The vertical line $x=0$ lies in $E$.

\emph{Count.} The number of squares in $\F$ (including $0$) is $(q+1)/2$: the squaring map is
$2$-to-$1$ on $\F^*$ because $-1\ne1$. Let $S_0$ be this set of squares. For each fixed $x\in\F$,
the map $s\mapsto x^2-s$ is a bijection from $S_0$ to $\{y: x^2-y\text{ is a square}\}$, so each
fixed $x$ contributes $(q+1)/2$ points. The first set therefore has $q(q+1)/2$ points. The vertical
line $\{(0,y)\}$ has $q$ points, of which $(q+1)/2$ already lie in the first set (those $y$ with $-y$
a square). The new points number $q-(q+1)/2=(q-1)/2$. Hence
$|E|=q(q+1)/2+(q-1)/2$.
\end{proof}

\begin{proof}[Proof of the Theorem]
Parts (a) and (b) give the minimum $\dfrac{q(q+1)}{2}+\dfrac{q-1}{2}$.
\end{proof}

\section{The Lean formalization}

The formalization is the single file \texttt{Kakeya.lean}. It imports only Mathlib, works for an
arbitrary finite field \texttt{F} with \texttt{[Field F] [Fintype F] [DecidableEq F]} (so for all odd
prime powers, not only primes), and contains no \texttt{sorry}, \texttt{native\_decide}, \texttt{axiom},
\texttt{unsafe}, \texttt{opaque}, \texttt{implemented\_by}, \texttt{extern} or \texttt{admit}. Declarations are in
the namespace \texttt{Kakeya}. The definitions below are quoted verbatim from the file, with the
Unicode symbols rendered in math mode.

\subsection{Definitions}

Lines (\texttt{lineS} is the non-vertical line $y=mx+b$; \texttt{lineV} the vertical line $x=c$):
\begin{lstlisting}
/-- The affine line `y = m x + b` (slope `m`, intercept `b`). -/
def lineS (m b : F) : Finset (F @$\times$@ F) := univ.image fun x => (x, m * x + b)
\end{lstlisting}
\begin{lstlisting}
/-- The vertical affine line `x = c`. -/
def lineV (c : F) : Finset (F @$\times$@ F) := univ.image fun y => (c, y)
\end{lstlisting}

Kakeya sets:
\begin{lstlisting}
/-- A Kakeya set: it contains a line `y = m x + b` for every slope `m`, and a vertical line. -/
def IsKakeya (K : Finset (F @$\times$@ F)) : Prop :=
  (@$\forall$@ m : F, @$\exists$@ b : F, @$\forall$@ x : F, (x, m * x + b) @$\in$@ K) @$\land$@ @$\exists$@ c : F, @$\forall$@ y : F, (c, y) @$\in$@ K
\end{lstlisting}

Line systems, multiplicity and cover:
\begin{lstlisting}
/-- The line of direction `d` in the system `t`. -/
def sysLine (t : Option F @$\to$@ F) (d : Option F) : Finset (F @$\times$@ F) := dirLine d (t d)
\end{lstlisting}
\begin{lstlisting}
/-- The multiplicity of a point: the number of lines of the system through it. -/
def mult (t : Option F @$\to$@ F) (A : F @$\times$@ F) : @$\mathbb{N}$@ :=
  (univ.filter fun d : Option F => A @$\in$@ sysLine t d).card
\end{lstlisting}
\begin{lstlisting}
/-- The union of the lines of the system, as the set of points of positive multiplicity. -/
def cover (t : Option F @$\to$@ F) : Finset (F @$\times$@ F) :=
  univ.filter fun A => 1 @$\le$@ mult t A
\end{lstlisting}

Excess and blocking sets:
\begin{lstlisting}
/-- The excess `2@$\sigma$@ = @$\sum$@_{mult A @$\ge$@ 1} (mult A - 1)(mult A - 2)`. -/
def excess (t : Option F @$\to$@ F) : @$\mathbb{Z}$@ :=
  @$\sum$@ A : F @$\times$@ F, if 1 @$\le$@ mult t A then ((mult t A : @$\mathbb{Z}$@) - 1) * ((mult t A : @$\mathbb{Z}$@) - 2) else 0
\end{lstlisting}
\begin{lstlisting}
/-- A blocking set: a set of points meeting every affine line. -/
def IsBlocking (B : Finset (F @$\times$@ F)) : Prop :=
  (@$\forall$@ m b : F, @$\exists$@ p @$\in$@ B, p.2 = m * p.1 + b) @$\land$@ @$\forall$@ c : F, @$\exists$@ p @$\in$@ B, p.1 = c
\end{lstlisting}
\begin{lstlisting}
/-- The polynomial `@$\prod$@_{(a,b) @$\in$@ B} (1 - a X@$_0$@ + b X@$_1$@)`. -/
noncomputable def blockPoly (B : Finset (F @$\times$@ F)) : MvPolynomial (Fin 2) F :=
  @$\prod$@ p @$\in$@ B, (1 - C p.1 * X 0 + C p.2 * X 1)
\end{lstlisting}

The extremal set:
\begin{lstlisting}
/-- Points `(x,y)` with `x@$^2$@ - y` a square: the union of the tangent lines `y = m x - m@$^2$@/4`
of the parabola `y = x@$^2$@` (the affine part of a dual conic). -/
def parabolaSet (F : Type*) [Field F] [Fintype F] [DecidableEq F] : Finset (F @$\times$@ F) :=
  univ.filter fun p => IsSquare (p.1 ^ 2 - p.2)
\end{lstlisting}
\begin{lstlisting}
/-- The extremal Kakeya set: `parabolaSet @$\cup$@ {x = 0}`. -/
def extremalSet (F : Type*) [Field F] [Fintype F] [DecidableEq F] : Finset (F @$\times$@ F) :=
  parabolaSet F @$\cup$@ lineV 0
\end{lstlisting}

\subsection{Main theorems}

The two main statements, quoted from the file:
\begin{lstlisting}
/-- **Blokhuis@$-$@Mazzocca.** For `q = |F|` odd, every Kakeya set in `F@$^2$@` has at least
`q(q+1)/2 + (q-1)/2` points, and this bound is attained. -/
theorem kakeya_min (hq : Odd (Fintype.card F)) :
    (@$\forall$@ K : Finset (F @$\times$@ F), IsKakeya K @$\to$@
        Fintype.card F * (Fintype.card F + 1) / 2 + (Fintype.card F - 1) / 2 @$\le$@ K.card) @$\land$@
      @$\exists$@ K : Finset (F @$\times$@ F), IsKakeya K @$\land$@
        K.card = Fintype.card F * (Fintype.card F + 1) / 2 + (Fintype.card F - 1) / 2 :=
\end{lstlisting}
\begin{lstlisting}
/-- The minimum size of a Kakeya set, as an infimum, equals `q(q+1)/2 + (q-1)/2` for `q` odd. -/
theorem kakeya_sInf (hq : Odd (Fintype.card F)) :
    sInf {n : @$\mathbb{N}$@ | @$\exists$@ K : Finset (F @$\times$@ F), IsKakeya K @$\land$@ K.card = n} =
      Fintype.card F * (Fintype.card F + 1) / 2 + (Fintype.card F - 1) / 2 := by
\end{lstlisting}

The supporting theorems are \texttt{card\_ge\_easy} (Faber's bound for every $q$),
\texttt{jamison} (blocking sets of $AG(2,q)$ have at least $2q-1$ points), \texttt{not\_allTwo\_pair}
(the product lemma), \texttt{card\_ge\_odd} (the refined bound for odd $q$), \texttt{isKakeya\_extremalSet}
and \texttt{card\_extremalSet} (the construction).

\subsection{Correspondence between the paper proof and the Lean declarations}

{\footnotesize\sloppy
\begin{tabular}{@{}p{3.6cm}p{5.6cm}p{6.2cm}@{}}
\toprule
Step of this report & Content & Lean declaration(s) \\
\midrule
3.1 line systems & $L_d$, union, Kakeya set contains a line system &
\texttt{lineS}, \texttt{lineV}, \texttt{sysLine}, \texttt{cover}, \texttt{cover\_subset} \\
3.1 intersections & distinct directions meet once; parallels disjoint &
\texttt{card\_inter\_dirLine}, \texttt{disjoint\_dirLine} \\
3.1 (C1), (C2) & double counts of Lemma~\ref{lem:counts} &
\texttt{sum\_mult}, \texttt{sum\_mult\_mul\_pred} \\
3.1 Faber's identity & $2|U|=q(q+1)+2\sigma$ &
\texttt{faber\_pointwise}, \texttt{cast\_mult\_mul\_pred}, \texttt{excess}, \texttt{excess\_nonneg}, \texttt{two\_mul\_card\_cover} \\
3.2 easy bound & $|K|\ge q(q+1)/2$ for every $q$ & \texttt{card\_ge\_easy} \\
3.3 Jamison & blocking polynomial, evaluation, degree sum & \texttt{IsBlocking}, \texttt{isBlocking\_translate}, \texttt{blockPoly}, \texttt{totalDegree\_blockPoly}, \texttt{eval\_blockPoly}, \texttt{eval\_blockPoly\_zero}, \texttt{eval\_blockPoly\_eq\_zero}, \texttt{sum\_eval\_eq\_zero}, \texttt{jamison} \\
3.4 (b) parity of lines & odd count on lines; all-2 dichotomy &
\texttt{sum\_mult\_line}, \texttt{sum\_mult\_line\_not\_sys}, \texttt{sum\_mult\_sysLine}, \texttt{exists\_odd\_mult\_of\_sum\_odd}, \texttt{exists\_odd\_not\_sys}, \texttt{AllTwo}, \texttt{exists\_odd\_or\_allTwo} \\
3.4 (c) product lemma & at most one non-vertical all-2 line &
\texttt{three\_le\_mult}, \texttt{allTwo\_not\_three}, \texttt{allTwo\_image}, \texttt{prod\_sub\_erase}, \texttt{not\_allTwo\_pair} \\
3.4 (a) $\tau\le4\sigma$ & pointwise inequality, summation &
\texttt{odd\_pointwise}, \texttt{card\_odd\_le} \\
3.4 (d) blocking set & point $E$, blocking property &
\texttt{exists\_blocking\_point}, \texttt{isBlocking\_odd} \\
3.4 (f) conclusion & $2\sigma\ge q-1$; refined bound &
\texttt{excess\_ge}, \texttt{card\_ge\_odd} \\
3.5 construction & parabola set, Kakeya property, count &
\texttt{parabolaSet}, \texttt{extremalSet}, \texttt{isKakeya\_extremalSet}, \texttt{two\_mul\_card\_squares}, \texttt{card\_squares}, \texttt{card\_fiber}, \texttt{card\_parabolaSet}, \texttt{card\_parabolaSet\_inter\_lineV}, \texttt{card\_extremalSet} \\
Theorem~\ref{thm:main} & both bounds and the infimum &
\texttt{kakeya\_min}, \texttt{kakeya\_sInf} \\
\bottomrule
\end{tabular}}

\subsection{Axioms and compilation status}

The file was recompiled for this report with \texttt{lake env lean} under Lean~4 toolchain
\texttt{v4.33.1} and Mathlib \texttt{v4.33.1}. It produced no errors and no warnings. The
\texttt{\#print axioms} commands at the end of the file report, for each of the eight declarations
\texttt{card\_ge\_easy}, \texttt{jamison}, \texttt{not\_allTwo\_pair}, \texttt{card\_ge\_odd},
\texttt{isKakeya\_extremalSet}, \texttt{card\_extremalSet}, \texttt{kakeya\_min}, \texttt{kakeya\_sInf}, exactly

\begin{center}
\texttt{depends on axioms: [propext, Classical.choice, Quot.sound]}
\end{center}

\noindent
which are the standard axioms of classical mathematics (no extra axiom is introduced).

A brute-force sanity check of the claimed minima is recorded separately, not in the Lean file. For
primes $p=3,5,7,11$ an exhaustive search (definition-level over all point sets for $p=3$,
over all line choices for $p=5,7$) and a branch-and-bound search (for $p=11$) give minima
$7,17,31,71$, matching the formula. For $q=13$ the formula gives $97$; the extremal set has $97$ points
by Section~3.5, which the Lean file proves for all odd $q$.

\section{Reproduction}

From the submission folder, the Lean part is rebuilt with
\begin{lstlisting}
cd lean && lake exe cache get && lake build
\end{lstlisting}
where the Lake project \texttt{lean/} contains the file \texttt{Kakeya.lean} in its \texttt{Submission}
library. The optional brute-force check in the submission folder is
\begin{lstlisting}
python3 check.py
\end{lstlisting}
which enumerates the small cases $p=3,5,7$ exhaustively and $p=11$ by branch and bound, and verifies
the minima $7,17,31,71$ against the formula $p(p+1)/2+(p-1)/2$. It is independent of the Lean proof and
only supports the sanity values above.

\end{document}
